% GENERATED FILE. Edit canonical lessons, then run npm run generate:books.
% canonical-source-hash: fnv1a64:37c3d3f757a5ed91
% canonical-lessons: SA-C494-lalayati, SA-C494-samvardhayati, SA-C494-bhartsayati, SA-C494-anujanati, SA-C494-balyam, SA-R494-review-488-494

\chapter{Bringing Up a Child}
\label{ch:sa-bringing-up-a-child}
\hlchaptermodality{494}

\begin{chapteropening}
\textbf{By the end of this chapter:} \emph{I can say how a child is brought up.}

Complete the last lesson of chapter 494: I can say how a child is brought up.
\end{chapteropening}

\section[\texorpdfstring{\sk{लालयति} (lālayati)}{लालयति (lālayati)}]{\sk{लालयति} (lālayati) — pampers}

\label{lesson:SA-C494-lalayati}



\begin{quote}
Before the new one: say the Sanskrit for the first haircut, then the Sanskrit for the sacred-thread ceremony.
\end{quote}

\subsection*{You'll want to know: \sk{लालयति}}
\textbf{\sk{लालयति}} — \emph{lālayati} — \textquotedblleft{}pampers\textquotedblright{}.

\begin{quote}
\sk{पितामही} \sk{पौत्रं} \sk{लालयति।}
\end{quote}

\begin{quote}
\emph{pitāmahī pautraṃ lālayati.}
\end{quote}

\begin{quote}
Grandmother pampers her grandson.
\end{quote}

Past (she …): \sk{सा} \sk{पौत्रं} \sk{लालितवती}, \emph{sā pautraṃ lālitavatī}.

\begin{cousinweb}
\textbf{\sk{लालयति}} (\emph{lālayati}) is built on the causative of the root √lal, \textquotedblleft{}to play, to frolic\textquotedblright{}; MW's ppp \emph{lālita} is \textquotedblleft{}caressed, fondled, indulged\textquotedblright{} (MW).
\end{cousinweb}

\begin{grammarlens}[title={What you've built}]
One more word for describing what you lived and felt.
\end{grammarlens}

\subsection*{Your turn}
\begin{itemize}
\raggedright
  \item \emph{Say it:} \emph{lālayati}
  \item \emph{Say it:} \emph{lālayati}, once more
  \item \emph{Say it:} \emph{pitāmahī pautraṃ lālayati.} — Grandmother pampers her grandson.
  \item \emph{Recall:} say the Sanskrit for the first haircut, then the Sanskrit for the sacred-thread ceremony, then say \emph{lālayati} again
\end{itemize}

\begin{tcolorbox}[breakable,colback=teal!4,colframe=teal!35!black,title={Before you move on}]
What does \sk{लालयति} mean? (\textquotedblleft{}Pampers\textquotedblright{}.) Say it once more.
\end{tcolorbox}
\section[\texorpdfstring{\sk{संवर्धयति} (saṃvardhayati)}{संवर्धयति (saṃvardhayati)}]{\sk{संवर्धयति} (saṃvardhayati) — brings up}

\label{lesson:SA-C494-samvardhayati}



\begin{quote}
Before the new one: say the Sanskrit for the sacred-thread ceremony, then the Sanskrit for pampers.
\end{quote}

\subsection*{You'll want to know: \sk{संवर्धयति}}
\textbf{\sk{संवर्धयति}} — \emph{saṃvardhayati} — \textquotedblleft{}brings up\textquotedblright{}.

\begin{quote}
\sk{माता} \sk{पिता} \sk{च} \sk{पुत्रं} \sk{स्नेहेन} \sk{संवर्धयतः।}
\end{quote}

\begin{quote}
\emph{mātā pitā ca putraṃ snehena saṃvardhayataḥ.}
\end{quote}

\begin{quote}
The mother and father bring up their son with love.
\end{quote}

Past (she …): \sk{सा} \sk{पुत्रीं} \sk{स्नेहेन} \sk{संवर्धितवती}, \emph{sā putrīṃ snehena saṃvardhitavatī}.

\begin{cousinweb}
\textbf{\sk{संवर्धयति}} (\emph{saṃvardhayati}) is built on \emph{sam} and the causative of the root √vṛdh, \textquotedblleft{}to grow\textquotedblright{}, the root of \textbf{\sk{वर्धते}}; MW's ppp \emph{saṃvardhita} is \textquotedblleft{}brought up, reared\textquotedblright{} (MW).
\end{cousinweb}

\begin{grammarlens}[title={What you've built}]
One more word for describing what you lived and felt.
\end{grammarlens}

\subsection*{Your turn}
\begin{itemize}
\raggedright
  \item \emph{Say it:} \emph{saṃvardhayati}
  \item \emph{Say it:} \emph{saṃvardhayati}, once more
  \item \emph{Say it:} \emph{mātā pitā ca putraṃ snehena saṃvardhayataḥ.} — The mother and father bring up their son with love.
  \item \emph{Recall:} say the Sanskrit for the sacred-thread ceremony, then the Sanskrit for pampers, then say \emph{saṃvardhayati} again
\end{itemize}

\begin{tcolorbox}[breakable,colback=teal!4,colframe=teal!35!black,title={Before you move on}]
What does \sk{संवर्धयति} mean? (\textquotedblleft{}Brings up\textquotedblright{}.) Say it once more.
\end{tcolorbox}
\section[\texorpdfstring{\sk{भर्त्सयति} (bhartsayati)}{भर्त्सयति (bhartsayati)}]{\sk{भर्त्सयति} (bhartsayati) — scolds}

\label{lesson:SA-C494-bhartsayati}



\begin{quote}
Before the new one: say the Sanskrit for pampers, then the Sanskrit for brings up.
\end{quote}

\subsection*{You'll want to know: \sk{भर्त्सयति}}
\textbf{\sk{भर्त्सयति}} — \emph{bhartsayati} — \textquotedblleft{}scolds\textquotedblright{}.

\begin{quote}
\sk{पिता} \sk{अलसं} \sk{पुत्रं} \sk{भर्त्सयति।}
\end{quote}

\begin{quote}
\emph{pitā alasaṃ putraṃ bhartsayati.}
\end{quote}

\begin{quote}
Father scolds his lazy son.
\end{quote}

Past (he …): \sk{सः} \sk{पुत्रं} \sk{भर्त्सितवान्}, \emph{saḥ putraṃ bhartsitavān}.

\begin{cousinweb}
\textbf{\sk{भर्त्सयति}} (\emph{bhartsayati}) comes from the root √bharts, \textquotedblleft{}to menace, to threaten, to abuse\textquotedblright{}; MW's ppp is \emph{bhartsita} (MW).
\end{cousinweb}

\begin{grammarlens}[title={What you've built}]
One more word for describing what you lived and felt.
\end{grammarlens}

\subsection*{Your turn}
\begin{itemize}
\raggedright
  \item \emph{Say it:} \emph{bhartsayati}
  \item \emph{Say it:} \emph{bhartsayati}, once more
  \item \emph{Say it:} \emph{pitā alasaṃ putraṃ bhartsayati.} — Father scolds his lazy son.
  \item \emph{Recall:} say the Sanskrit for pampers, then the Sanskrit for brings up, then say \emph{bhartsayati} again
\end{itemize}

\begin{tcolorbox}[breakable,colback=teal!4,colframe=teal!35!black,title={Before you move on}]
What does \sk{भर्त्सयति} mean? (\textquotedblleft{}Scolds\textquotedblright{}.) Say it once more.
\end{tcolorbox}
\section[\texorpdfstring{\sk{अनुजानाति} (anujānāti)}{अनुजानाति (anujānāti)}]{\sk{अनुजानाति} (anujānāti) — permits}

\label{lesson:SA-C494-anujanati}



\begin{quote}
Before the new one: say the Sanskrit for brings up, then the Sanskrit for scolds.
\end{quote}

\subsection*{You'll want to know: \sk{अनुजानाति}}
\textbf{\sk{अनुजानाति}} — \emph{anujānāti} — \textquotedblleft{}permits\textquotedblright{}.

\begin{quote}
\sk{पिता} \sk{पुत्रस्य} \sk{यात्राम्} \sk{अनुजानाति।}
\end{quote}

\begin{quote}
\emph{pitā putrasya yātrām anujānāti.}
\end{quote}

\begin{quote}
Father allows his son's journey.
\end{quote}

Past (he …): \sk{सः} \sk{यात्राम्} \sk{अनुज्ञातवान्}, \emph{saḥ yātrām anujñātavān}.

\begin{cousinweb}
\textbf{\sk{अनुजानाति}} (\emph{anujānāti}) is built from \emph{anu}, \textquotedblleft{}after\textquotedblright{}, and the root √jñā, \textquotedblleft{}to know\textquotedblright{}, the root of \textbf{\sk{जानाति}}; MW gives \textquotedblleft{}to permit, grant, allow\textquotedblright{} (MW).
\end{cousinweb}

\begin{grammarlens}[title={What you've built}]
One more word for describing what you lived and felt.
\end{grammarlens}

\subsection*{Your turn}
\begin{itemize}
\raggedright
  \item \emph{Say it:} \emph{anujānāti}
  \item \emph{Say it:} \emph{anujānāti}, once more
  \item \emph{Say it:} \emph{pitā putrasya yātrām anujānāti.} — Father allows his son's journey.
  \item \emph{Recall:} say the Sanskrit for brings up, then the Sanskrit for scolds, then say \emph{anujānāti} again
\end{itemize}

\begin{tcolorbox}[breakable,colback=teal!4,colframe=teal!35!black,title={Before you move on}]
What does \sk{अनुजानाति} mean? (\textquotedblleft{}Permits\textquotedblright{}.) Say it once more.
\end{tcolorbox}
\section[\texorpdfstring{\sk{बाल्यम्} (bālyam)}{बाल्यम् (bālyam)}]{\sk{बाल्यम्} (bālyam) — childhood}

\label{lesson:SA-C494-balyam}



\begin{quote}
Before the new one: say the Sanskrit for scolds, then the Sanskrit for permits.
\end{quote}

\subsection*{You'll want to know: \sk{बाल्यम्}}
\textbf{\sk{बाल्यम्}} — \emph{bālyam} — \textquotedblleft{}childhood\textquotedblright{}.

\begin{quote}
\sk{बाल्ये} \sk{अहं} \sk{ग्रामे} \sk{अवसम्।}
\end{quote}

\begin{quote}
\emph{bālye ahaṃ grāme avasam.}
\end{quote}

\begin{quote}
In childhood I lived in the village.
\end{quote}

A neuter noun.

\begin{cousinweb}
\textbf{\sk{बाल्यम्}} (\emph{bālyam}) is built on \textbf{\sk{बालः}}, \textquotedblleft{}a child\textquotedblright{}, by lengthening its vowel and adding \emph{-ya}; MW gives \textquotedblleft{}boyhood, childhood, infancy\textquotedblright{} (MW).
\end{cousinweb}

\begin{grammarlens}[title={What you've built}]
One more word for describing what you lived and felt.
\end{grammarlens}

\subsection*{Your turn}
\begin{itemize}
\raggedright
  \item \emph{Say it:} \emph{bālyam}
  \item \emph{Say it:} \emph{bālyam}, once more
  \item \emph{Say it:} \emph{bālye ahaṃ grāme avasam.} — In childhood I lived in the village.
  \item \emph{Recall:} say the Sanskrit for scolds, then the Sanskrit for permits, then say \emph{bālyam} again
\end{itemize}

\begin{tcolorbox}[breakable,colback=teal!4,colframe=teal!35!black,title={Before you move on}]
What does \sk{बाल्यम्} mean? (\textquotedblleft{}Childhood\textquotedblright{}.) Say it once more.
\end{tcolorbox}
\section[(dialogue)]{Looking back: from Sleep and Waking to Bringing Up a Child}

\label{lesson:SA-R494-review-488-494}



\begin{quote}
Say the Sanskrit for awakes and lassitude.
\end{quote}

\subsection*{Your turn}
Say these aloud:

\begin{itemize}
\raggedright
  \item awakes, lassitude, wakes someone up, sleepless, a three-hour watch
  \item a vice, self-control, smoking, gives up entirely, liquor
  \item pale, length of life, grey-haired, a wrinkle, appearance
  \item old age, wears out, death, funeral rites, a rite for the ancestors
  \item a pregnant woman, gives birth, birth, newborn, a swing
  \item a rite of passage, the naming ceremony, a baby's first rice, the first haircut, the sacred-thread ceremony
  \item pampers, brings up, scolds, permits, childhood
\end{itemize}

\begin{tcolorbox}[breakable,colback=teal!4,colframe=teal!35!black,title={Before you move on}]
Awakes? (\textbf{\sk{प्रतिबुध्यते}}, \emph{pratibudhyate}.) Lassitude? (\textbf{\sk{तन्द्रा}}, \emph{tandrā}.) Wakes someone up? (\textbf{\sk{प्रबोधयति}}, \emph{prabodhayati}.) Sleepless? (\textbf{\sk{उन्निद्रः}}, \emph{unnidraḥ}.) A three-hour watch? (\textbf{\sk{प्रहरः}}, \emph{praharaḥ}.) A vice? (\textbf{\sk{व्यसनम्}}, \emph{vyasanam}.) Self-control? (\textbf{\sk{संयमः}}, \emph{saṃyamaḥ}.) Smoking? (\textbf{\sk{धूमपानम्}}, \emph{dhūmapānam}.) Gives up entirely? (\textbf{\sk{परित्यजति}}, \emph{parityajati}.) Liquor? (\textbf{\sk{मद्यम्}}, \emph{madyam}.) Pale? (\textbf{\sk{गौरः}}, \emph{gauraḥ}.) Length of life? (\textbf{\sk{आयुः}}, \emph{āyuḥ}.) Grey-haired? (\textbf{\sk{पलितः}}, \emph{palitaḥ}.) A wrinkle? (\textbf{\sk{वलिः}}, \emph{valiḥ}.) Appearance? (\textbf{\sk{रूपम्}}, \emph{rūpam}.) Old age? (\textbf{\sk{जरा}}, \emph{jarā}.) Wears out? (\textbf{\sk{जीर्यति}}, \emph{jīryati}.) Death? (\textbf{\sk{मृत्युः}}, \emph{mṛtyuḥ}.) Funeral rites? (\textbf{\sk{अन्त्येष्टिः}}, \emph{antyeṣṭiḥ}.) A rite for the ancestors? (\textbf{\sk{श्राद्धम्}}, \emph{śrāddham}.) A pregnant woman? (\textbf{\sk{गर्भिणी}}, \emph{garbhiṇī}.) Gives birth? (\textbf{\sk{प्रसूते}}, \emph{prasūte}.) Birth? (\textbf{\sk{जन्म}}, \emph{janma}.) Newborn? (\textbf{\sk{नवजातः}}, \emph{navajātaḥ}.) A swing? (\textbf{\sk{दोला}}, \emph{dolā}.) A rite of passage? (\textbf{\sk{संस्कारः}}, \emph{saṃskāraḥ}.) The naming ceremony? (\textbf{\sk{नामकरणम्}}, \emph{nāmakaraṇam}.) A baby's first rice? (\textbf{\sk{अन्नप्राशनम्}}, \emph{annaprāśanam}.) The first haircut? (\textbf{\sk{चूडाकरणम्}}, \emph{cūḍākaraṇam}.) The sacred-thread ceremony? (\textbf{\sk{उपनयनम्}}, \emph{upanayanam}.) Pampers? (\textbf{\sk{लालयति}}, \emph{lālayati}.) Brings up? (\textbf{\sk{संवर्धयति}}, \emph{saṃvardhayati}.) Scolds? (\textbf{\sk{भर्त्सयति}}, \emph{bhartsayati}.) Permits? (\textbf{\sk{अनुजानाति}}, \emph{anujānāti}.) Childhood? (\textbf{\sk{बाल्यम्}}, \emph{bālyam}.)
\end{tcolorbox}
